% !TeX root = ../教材题库_全部条目_转写工作稿.tex
% Source: XJ-B1-C06-PR-P241-Y1232
% 第6章 / 6.4 用样本估计总体 / 教材 p.241 / PDF 第248页
% OCR未校对：公式与文字均可能有误；未完成数学TeX转写；禁止直接发给学生
% 原题及相邻原解截图（核对用）：../题目原图/XJ-B1-C06-PR-P241-Y1232-01.webp
% 本文件只在转写完成后作为题干使用；OCR 条目尚不是可靠题干。
\begin{bankitem}{p.241 练习 第2题}{第6章 / 6.4 用样本估计总体 / 教材 p.241 / PDF 第248页}{XJ-B1-C06-PR-P241-Y1232}{OCR未校对：公式与文字均可能有误；未完成数学TeX转写；禁止直接发给学生}
\begin{ocrdraft}
[教材p.241 / PDF第248页]
\par
2．甲、乙、丙三家电子厂商都在广告中声称，他们的某型电子产品在正常
\par
情况下的待机时间是12h，质量检测部门对这三家销售产品的待机时间进行了抽
\par
样调查，统计结果（单位：h）如下：
\par
甲：
\par
8，9，9，9，9，11，13，16，17，19；
\par
乙：
\par
10，10，12，12，12，13，14，16，18，19；
\par
丙：
\par
8，8，8，10，11，
\par
13，17，19，20，20．
\par
分别求出以上三组数据的平均数、众数、中位数．
\par
（2）这三个厂商的推销广告分别利用了上述哪一种数据来表示待机时间？
\par
（3）如果你是顾客，宜选择哪个厂商的产品？为什么？
\end{ocrdraft}
\end{bankitem}
